\section{Extended results}\label{app:extra}

\paragraph{Accuracy by target and evaluation scheme.} Table~\ref{tab:accuracy} gives all three targets at the four levels,
with settings chosen inside each split (nested) and with the deployed settings. Table~\ref{tab:pressure} gives the
baselines with bootstrap intervals and the shuffled-label control.

\begin{table}[ht]\centering\small
\caption{Accuracy by closest labeled relative, mean of 3 seeds. Error is the typical fold error for doubling times and the mean absolute error in \textdegree C for growth temperature; $\rho$ is the Spearman correlation. Nested: settings chosen inside each split, each training genome entered once per placement rank, as deployed (the estimates quoted in the main text). Fixed, coarsened: the deployed settings, chosen on all rows, with the deployed training scheme. Fixed, single: the deployed settings and each genome entered once.}\label{tab:accuracy}
\begin{tabular}{lcccccccc}
\toprule
& \multicolumn{2}{c}{same genus} & \multicolumn{2}{c}{new genus} & \multicolumn{2}{c}{new family} & \multicolumn{2}{c}{new order} \\
\cmidrule(lr){2-3} \cmidrule(lr){4-5} \cmidrule(lr){6-7} \cmidrule(lr){8-9}
& Error & $\rho$ & Error & $\rho$ & Error & $\rho$ & Error & $\rho$ \\
\midrule
\multicolumn{9}{l}{\textit{Doubling time at the growth temperature}} \\
\quad nested & 1.75 & 0.74 & 1.91 & 0.71 & 2.06 & 0.66 & 2.08 & 0.66 \\
\quad fixed, coarsened & 1.75 & 0.74 & 1.92 & 0.71 & 2.07 & 0.67 & 2.08 & 0.66 \\
\quad fixed, single & 1.82 & 0.74 & 1.96 & 0.71 & 2.07 & 0.67 & 2.08 & 0.67 \\
\multicolumn{9}{l}{\textit{Doubling time at 20 \textdegree C}} \\
\quad nested & 1.73 & 0.84 & 1.81 & 0.83 & 1.91 & 0.81 & 1.93 & 0.81 \\
\quad fixed, coarsened & 1.73 & 0.84 & 1.81 & 0.83 & 1.88 & 0.82 & 1.93 & 0.81 \\
\quad fixed, single & 1.75 & 0.84 & 1.79 & 0.83 & 1.88 & 0.82 & 1.91 & 0.81 \\
\multicolumn{9}{l}{\textit{Growth temperature}} \\
\quad nested & 4.96 & 0.86 & 5.74 & 0.83 & 5.81 & 0.83 & 5.99 & 0.81 \\
\quad fixed, coarsened & 5.13 & 0.85 & 5.77 & 0.83 & 5.87 & 0.82 & 6.07 & 0.81 \\
\quad fixed, single & 5.36 & 0.85 & 5.88 & 0.82 & 5.96 & 0.82 & 6.15 & 0.81 \\
\bottomrule
\end{tabular}
\end{table}

\begin{table}[ht]\centering\small
\caption{Pressure tests of the headline (doubling time at the growth temperature, mean of 3 seeds, the folds of the headline). Intervals: 95\% bootstrap over species groups; for this model, of the typical error and Spearman correlation; for the median of the closest relatives, of the model's error minus the baseline's (negative: the model is better). Shuffled labels: the deployed model trained and scored on labels permuted across genomes.}\label{tab:pressure}
\begin{tabular}{lcccc}
\toprule
& Error ($\times$) & Spearman & AUC, slow & 95\% interval \\
\midrule
\multicolumn{5}{l}{\textit{Same genus}} \\
\quad this model, nested & 1.75 & 0.74 & 0.85 & 1.67--1.86, 0.68--0.78 \\
\quad median of closest relatives & 1.81 & 0.66 & 0.80 & $-$0.21 to +0.05 \\
\quad median training label & 2.72 & $-$0.08 & 0.45 &  \\
\multicolumn{5}{l}{\textit{New genus}} \\
\quad this model, nested & 1.91 & 0.71 & 0.84 & 1.81--2.02, 0.66--0.76 \\
\quad median of closest relatives & 2.12 & 0.57 & 0.79 & $-$0.42 to $-$0.02 \\
\quad median training label & 2.74 & $-$0.07 & 0.47 &  \\
\multicolumn{5}{l}{\textit{New family}} \\
\quad this model, nested & 2.06 & 0.66 & 0.81 & 1.93--2.24, 0.59--0.72 \\
\quad median of closest relatives & 2.59 & 0.40 & 0.72 & $-$0.84 to $-$0.21 \\
\quad median training label & 2.86 & $-$0.20 & 0.43 &  \\
\multicolumn{5}{l}{\textit{New order}} \\
\quad this model, nested & 2.08 & 0.66 & 0.82 & 1.96--2.23, 0.57--0.72 \\
\quad median of closest relatives & 2.45 & 0.41 & 0.72 & $-$0.74 to $-$0.16 \\
\quad median training label & 2.78 & $-$0.18 & 0.40 &  \\
\multicolumn{5}{l}{\textit{Labels shuffled across genomes}} \\
\quad this model, species held out & 2.91 & 0.00 & 0.49 & \\
\quad median training label, species held out & 2.73 & $-$0.03 & 0.50 & \\
\quad this model, order held out & 2.89 & $-$0.04 & 0.47 & \\
\quad median training label, order held out & 2.72 & $-$0.06 & 0.45 & \\
\bottomrule
\end{tabular}
\end{table}


\paragraph{Ranges and the slow-growth probability.} Table~\ref{tab:ranges} gives the coverage of the 80\% ranges by group,
cross-fitted over genera, and Table~\ref{tab:slowprob} the probability of a doubling time above a day.

\begin{table}[ht]\centering\small
\caption{Coverage (\%) of the 80\% ranges for the doubling time at the growth temperature, cross-fitted over genera, by group of genomes, for ranges set by novelty level alone and by novelty level and domain (released). Groups by the true doubling time cannot be targeted by any range built from the inputs.}\label{tab:ranges}
\begin{tabular}{llccccc}
\toprule
Held out & Ranges by & All & Bacteria & Archaea & Fast in truth & Slow in truth \\
\midrule
same genus & novelty & 80 & 80 & 79 & 83 & 44 \\
 & novelty, domain & 80 & 80 & 79 & 83 & 45 \\
new genus & novelty & 80 & 80 & 80 & 83 & 36 \\
 & novelty, domain & 79 & 79 & 79 & 83 & 37 \\
new family & novelty & 80 & 79 & 81 & 84 & 29 \\
 & novelty, domain & 80 & 80 & 80 & 84 & 29 \\
new order & novelty & 80 & 82 & 73 & 85 & 26 \\
 & novelty, domain & 79 & 79 & 78 & 83 & 29 \\
\bottomrule
\end{tabular}
\end{table}

\begin{table}[ht]\centering\small
\caption{The probability of a doubling time above 24 h (logistic model on the predicted log doubling time and the domain), cross-fitted over genera, by the closest labeled relative. Calibration error: mean absolute difference between predicted probability and observed share over fifths of the predicted probability.}\label{tab:slowprob}
\begin{tabular}{lcccc}
\toprule
& same genus & new genus & new family & new order \\
\midrule
AUC & 0.84 & 0.83 & 0.79 & 0.81 \\
Brier score & 0.057 & 0.060 & 0.066 & 0.065 \\
Brier score of the base rate & 0.072 & 0.072 & 0.072 & 0.072 \\
calibration error & 0.010 & 0.008 & 0.009 & 0.012 \\
slow among the top fifth & 0.24 & 0.25 & 0.22 & 0.22 \\
slow overall & 0.08 & 0.08 & 0.08 & 0.08 \\
\bottomrule
\end{tabular}
\end{table}


\paragraph{Ablations.} Table~\ref{tab:ablations} removes one input group at a time from the feature model (deployed
settings, each genome once, \NSeedsAblation seeds).

\begin{table}[ht]\centering\scriptsize\setlength{\tabcolsep}{1.5pt}
\caption{What each part contributes. Each row changes one thing in the final model and gives the mean change over 10 seeds, with the number of seeds in which the variant did better in brackets; $^{*}$ paired Wilcoxon $p<0.05$. Typical error is the fold error of the median genome (lower is better); AUC ranks genomes slower than a day. All variants are scored on the same genomes and folds, with the deployed settings and each genome entered once. The seeds reshuffle the same genomes, so $p$ measures how consistent a change is across splits, not independent replication. $^{\dagger}$The earlier learner (gradient boosting, 200 trees of depth 3) against the tuned average, both with settings chosen inside each split (5 seeds).}\label{tab:ablations}
\begin{tabular}{lcccccc}
\toprule
& \multicolumn{3}{c}{Species held out} & \multicolumn{3}{c}{Orders held out} \\
\cmidrule(lr){2-4} \cmidrule(lr){5-7}
& Typical error ($\times$) & Spearman & AUC, slow & Typical error ($\times$) & Spearman & AUC, slow \\
\midrule
final model & 1.819 & 0.740 & 0.838 & 2.093 & 0.654 & 0.803 \\
\midrule
without gRodon statistics & +0.110$^{*}$ (0/10) & $-$0.044$^{*}$ (0/10) & $-$0.010$^{*}$ (0/10) & +0.295$^{*}$ (0/10) & $-$0.080$^{*}$ (0/10) & $-$0.045$^{*}$ (0/10) \\
without lifestyle genes & +0.005 (4/10) & $-$0.002$^{*}$ (1/10) & 0.000 (3/10) & $-$0.020 (6/10) & +0.005$^{*}$ (8/10) & +0.002$^{*}$ (7/10) \\
without any gene features & +0.274$^{*}$ (0/10) & $-$0.078$^{*}$ (0/10) & $-$0.033$^{*}$ (0/10) & +0.551$^{*}$ (0/10) & $-$0.195$^{*}$ (0/10) & $-$0.076$^{*}$ (0/10) \\
without 16S & $-$0.002 (7/10) & +0.002$^{*}$ (8/10) & 0.000 (6/10) & 0.000 (6/10) & 0.000 (7/10) & $-$0.010$^{*}$ (1/10) \\
without GTDB phylogeny & +0.026$^{*}$ (1/10) & $-$0.008$^{*}$ (0/10) & +0.002 (7/10) & $-$0.023$^{*}$ (8/10) & +0.002 (7/10) & +0.006$^{*}$ (10/10) \\
without the added species & +0.040$^{*}$ (0/10) & $-$0.009$^{*}$ (0/10) & $-$0.003$^{*}$ (1/10) & +0.040$^{*}$ (0/10) & $-$0.004 (3/10) & $-$0.006$^{*}$ (2/10) \\
with energy-metabolism genes & +0.002 (4/10) & 0.000 (5/10) & +0.001 (7/10) & +0.013$^{*}$ (3/10) & $-$0.001 (4/10) & $-$0.001 (4/10) \\
with predicted growth temperature & $-$0.015 (6/10) & $-$0.001 (4/10) & 0.000 (5/10) & $-$0.015$^{*}$ (8/10) & +0.004$^{*}$ (8/10) & +0.002 (7/10) \\
earlier learner, nested$^{\dagger}$ & & & & +0.015 (1/5) & $-$0.021 (0/5) & $-$0.012 (0/5) \\
\bottomrule
\end{tabular}
\end{table}


\paragraph{Excluded labels.} Table~\ref{tab:exclusions} lists the removed labels with their evidence.

{\small
\begin{longtable}{p{2.4cm}p{1.0cm}p{5.2cm}p{3.5cm}}
\caption{Doubling times removed from the training data because they are not maximal growth rates, with the evidence.}\label{tab:exclusions}\\
\toprule
Species & Label (h) & Evidence & Source \\
\midrule
\endfirsthead
\caption[]{(continued)}\\
\toprule
Species & Label (h) & Evidence & Source \\
\midrule
\endhead
\textit{Methylocystis parvus} & 152 & source gives about 3700 h at 30 C for a methanotroph; same study gives Methylococcus capsulatus about 4600 h against 1.87 h elsewhere & Kevbrina et al. 2001 Microbiology 70:384 (growth of methanotrophs at low temperatures) via Corkrey et al. \\
\textit{Methylomonas methanica} & 112 & source gives about 2700 h; same study as above & Kevbrina et al. 2001 Microbiology 70:384 via Corkrey et al. \\
\textit{Methylosinus sporium} & 187 & source gives about 4600 h; same study as above & Kevbrina et al. 2001 Microbiology 70:384 via Corkrey et al. \\
\textit{Methylosinus trichosporium} & 131 & source gives about 3200 h; same study as above & Kevbrina et al. 2001 Microbiology 70:384 via Corkrey et al. \\
\textit{Flavobacterium segetis} & 433 & recorded at 12.9 C; Bergey's Manual gives 6.8 h; genus relatives median 4.9 h & Yi and Chun 2006 (species description) via Corkrey et al. \\
\textit{Flavobacterium weaverense} & 178 & recorded at 15.6 C; Bergey's Manual gives 2.9 h & Yi and Chun 2006 (species description) via Corkrey et al. \\
\textit{Acidithiobacillus ferrooxidans} & 386 & about 70 times its genus relatives (median 5.5 h); other organisms from the same study look plausible, cause unknown & Franzmann et al. (temperature and iron/sulfur oxidation rates) via Corkrey et al. \\
\textit{Acidithiobacillus thiooxidans} & 74 & about 13 times its genus relatives; same study as above & Franzmann et al. via Corkrey et al. \\
\textit{Klebsiella oxytoca} & 11.7 & from a lag-phase study in meat; genus relatives median 0.3 h & Mellefont (meat industry predictive modelling) via Corkrey et al. \\
\textit{Streptomyces coelicolor} & 0.5 & another source gives 2.04 h; MediaDB values are usually about 1.6 times slower than other sources (defined media), this one is 4 times faster, and 0.5 h is implausible for a filamentous Streptomyces & MediaDB (Richards et al. 2014) via Madin et al.; found by screen\_label\_sources.py \\
\bottomrule
\end{longtable}
}


\paragraph{External test set.} Table~\ref{tab:external} lists the external species with their measured and predicted
doubling times and the source of each value.

{\footnotesize\setlength{\tabcolsep}{3pt}
\begin{longtable}{llcccccl}
\caption{External test: 48 species whose doubling time comes from the species description and that are in neither the Madin compilation (by NCBI taxonomy identifier) nor our labels, with no GTDB species in training. Measured, predicted and relatives' median doubling times in hours at the growth temperature; $T$: the growth temperature stated with the value (\textdegree C, blank where none was stated); range: whether the measured value lies inside the released 80\% range.}\label{tab:external}\\
\toprule
Species & Closest labeled relative & $T$ & Measured & Predicted & Relatives & Range & Source \\
\midrule
\endfirsthead
\toprule
Species & Closest labeled relative & $T$ & Measured & Predicted & Relatives & Range & Source \\
\midrule
\endhead
\textit{Carboxydothermus islandicus} & same genus &  & 2 & 1.94 & 2 & yes & PMID 21131500 \\
\textit{Carboxydothermus pertinax} & same genus &  & 1.5 & 3.19 & 2 & yes & PMID 21908679 \\
\textit{Hydrogenimonas urashimensis} & same genus &  & 3 & 2.99 & 1.16 & yes & PMID 33340909 \\
\textit{Methanospirillum purgamenti} & same genus & 35 & 10 & 9.78 & 20.7 & yes & PMID 39186748 \\
\textit{Methylacidiphilum kamchatkense} & same genus & 55 & 6.3 & 8.73 & 9.9 & yes & PMID 37755432 \\
\textit{Paracoccus solventivorans} & same genus &  & 12 & 5.21 & 2.1 & yes & PMID 8863446 \\
\textit{Thermocrinis jamiesonii} & same genus &  & 1.3 & 1.5 & 1.18 & yes & PMID 26419502 \\
\textit{Aminomonas paucivorans} & same family, new genus & 35 & 16 & 2.36 & 63.3 & no & PMID 10425753 \\
\textit{Anaerobium acetethylicum} & same family, new genus & 30 & 1 & 0.59 & 4.35 & yes & PMID 26297346 \\
\textit{Anaeromicrobium sediminis} & same family, new genus & 37 & 1.27 & 1.7 & 5.73 & yes & PMID 27974087 \\
\textit{Aureliella helgolandensis} & same family, new genus &  & 41 & 15.2 & 10 & yes & PMID 32219667 \\
\textit{Caloramator quimbayensis} & same family, new genus & 50 & 2 & 2.63 & 0.324 & yes & PMID 22843719 \\
\textit{Methanothermococcus jasoni} & same family, new genus & 73 & 0.533 & 0.688 & 0.917 & yes & PMID 41773501 \\
\textit{Ornatilinea apprima} & same family, new genus &  & 6 & 48.7 & 50 & no & PMID 22328612 \\
\textit{Persephonella hydrogeniphila} & same family, new genus &  & 0.667 & 2.2 & 2.44 & yes & PMID 12807214 \\
\textit{Phorcysia thermohydrogeniphila} & same family, new genus &  & 0.433 & 1.15 & 1.56 & yes & PMID 22140151 \\
\textit{Spiroplasma corruscae} & same family, new genus & 32 & 1.5 & 4.4 & 2.27 & yes & PMID 8863421 \\
\textit{Spiroplasma eriocheiris} & same family, new genus &  & 24 & 5.11 & 2.27 & no & PMID 20418415 \\
\textit{Spiroplasma poulsonii} & same family, new genus & 30 & 15.8 & 9.82 & 2.27 & yes & PMID 10319483 \\
\textit{Spiroplasma syrphidicola} & same family, new genus &  & 1 & 2.55 & 2.27 & yes & PMID 8782692 \\
\textit{Sulfurovum aggregans} & same family, new genus &  & 3.2 & 1.64 & 1.84 & yes & PMID 24966202 \\
\textit{Sulfurovum riftiae} & same family, new genus &  & 3 & 1.61 & 1.84 & yes & PMID 27116914 \\
\textit{Tautonia plasticadhaerens} & same family, new genus &  & 29 & 8.54 & 17.4 & no & PMID 32399714 \\
\textit{Thermocrinis minervae} & same family, new genus & 75 & 4.5 & 1.97 & 1.4 & yes & PMID 19651724 \\
\textit{Varunaivibrio sulfuroxidans} & same family, new genus &  & 8 & 8.56 & 11 & yes & PMID 27298285 \\
\textit{Deferribacter autotrophicus} & same order, new family &  & 1 & 5.54 & 8.42 & no & PMID 19502344 \\
\textit{Deferribacter desulfuricans} & same order, new family &  & 0.667 & 6.19 & 8.42 & no & PMID 12807210 \\
\textit{Defluviitalea raffinosedens} & same order, new family & 50 & 1.2 & 3 & 4.35 & yes & PMC5817277 \\
\textit{Desulfovulcanus ferrireducens} & same order, new family & 65 & 7.1 & 11.1 & 16.7 & yes & PMID 35190935 \\
\textit{Fervidicoccus fontis} & same order, new family & 67.5 & 6.1 & 3.2 & 3.71 & yes & PMID 19837732 \\
\textit{Lentimicrobium saccharophilum} & same order, new family &  & 24 & 11.8 & 1.47 & yes & PMID 27098854 \\
\textit{Spiroplasma platyhelix} & same order, new family & 30 & 6.4 & 2.36 & 2.72 & yes & PMID 9226909 \\
\textit{Sporanaerobium hydrogeniformans} & same order, new family &  & 5 & 0.87 & 4.35 & no & PMID 37572166 \\
\textit{Sulfurimonas hongkongensis} & same order, new family & 30 & 6.1 & 1.73 & 1.22 & yes & PMC4149026 \\
\textit{Aggregatilinea lenta} & new order &  & 456 & 11.2 & 50 & no & PMID 30775966 \\
\textit{Anoxybacter fermentans} & new order & 60 & 0.45 & 1.37 & 3.5 & yes & PMID 25505345 \\
\textit{Atrimonas thermophila} & new order & 72.5 & 11.4 & 2.2 & 2.1 & no & PMID 38776610 \\
\textit{Brockia lithotrophica} & new order & 62.5 & 6 & 1.3 & 0.5 & no & PMID 22493174 \\
\textit{Caldanaerovirga acetigignens} & new order &  & 3.4 & 2.73 & 1.25 & yes & PMID 19625440 \\
\textit{Calderihabitans maritimus} & new order & 65 & 3.7 & 3.44 & 1.25 & yes & PMID 23606483 \\
\textit{Carboxydocella thermautotrophica} & new order &  & 1.1 & 1.54 & 1.25 & yes & PMID 12508854 \\
\textit{Deferrisoma palaeochoriense} & new order &  & 2.5 & 2.77 & 12.4 & yes & PMID 26610851 \\
\textit{Desulfobacca acetoxidans} & new order & 37 & 48 & 10.1 & 12.4 & no & PMID 10319454 \\
\textit{Desulfocarbo indianensis} & new order & 30 & 64.8 & 13.4 & 12.4 & no & PMID 24876241 \\
\textit{Mariprofundus micogutta} & new order & 25 & 24 & 2.48 & 1.73 & no & PMID 27766355 \\
\textit{Phosphitispora fastidiosa} & new order & 30 & 72 & 7.39 & 1.25 & no & PMID 34878375 \\
\textit{Stratiformator vulcanicus} & new order &  & 102 & 20.5 & 13.2 & no & PMID 37584762 \\
\textit{Thermincola carboxydiphila} & new order &  & 1.3 & 6 & 1.25 & no & PMID 16166711 \\
\bottomrule
\end{longtable}}


\section{Implementation and compute}\label{app:compute}
All experiments ran on one laptop (Apple silicon, 10 CPU cores, 16 GB memory). Gene prediction and the HMM searches for the
labeled genomes take under an hour on eight cores; ESM-2 (35M) embeddings of the sampled proteins of all labeled genomes ran
on the laptop GPU in about an hour; the nested evaluation at four levels and three seeds takes several hours on nine cores.
Software versions are listed in the repository. Splits are defined by the seed and the GTDB taxonomy of each genome, and the
released code regenerates them exactly.
